\section{Implementation Strategy}

The subtraction formalism is implemented at the level of the Monte Carlo
integrand, where each phase-space point contributes a finite weight.
Particular care is taken to follow closely the structure of modern
integrators such as the Munich framework used in MATRIX.

\subsection{Monte Carlo weight}

At each phase-space point, the weight is given by
\begin{align}
w =
\frac{1}{2\hat{s}} \,
f_a(x_1,\mu_F) \, f_b(x_2,\mu_F)
\, d\Phi \,
\alpha_s^n(\mu_R)
\, |M|^2,
\end{align}
where $\hat{s} = x_1 x_2 s$ and $d\Phi$ includes all Jacobians from the
phase-space mappings.

\subsection{Local subtraction at the integrand level}

The real-emission contribution is evaluated as
\begin{align}
\mathcal{R}_{\text{finite}} =
|M_{m+1}|^2 - \sum_{\text{dipoles}} D_{ij,k},
\end{align}
which is finite at every phase-space point.

The virtual contribution is combined with the integrated dipoles:
\begin{align}
\mathcal{V}_{\text{finite}} =
|M_m^{\text{1-loop}}|^2 + \langle M_m | \mathbf{I} | M_m \rangle.
\end{align}

\subsection{Multichannel phase-space sampling}

Following the strategy used in Munich, the phase space is sampled using
a multichannel approach. The idea is to construct a probability density
as a weighted sum over different phase-space channels:
\begin{align}
g_{\text{MC}}(x) = \sum_c \alpha_c \, g_c(x),
\end{align}
where:
\begin{itemize}
    \item $g_c(x)$ are individual channel densities adapted to specific
    kinematic structures (e.g. soft or collinear regions),
    \item $\alpha_c$ are channel weights satisfying $\sum_c \alpha_c = 1$.
\end{itemize}

Each channel is designed to capture a particular singular behavior of
the matrix element, improving the efficiency of the integration.

\subsection{Adaptive integration (VEGAS)}

On top of the multichannel structure, an adaptive VEGAS algorithm is used
to refine the sampling of the integration variables.

The VEGAS procedure builds an importance-sampling grid for each integration
dimension, iteratively updating the sampling density according to
\begin{align}
g(x) \propto |w(x)|.
\end{align}

This allows the integrator to concentrate sampling in regions where the
integrand is large, such as near infrared-sensitive configurations.

\subsection{Combination of multichannel and VEGAS}

The final sampling density is given by the combination of both strategies:
\begin{itemize}
    \item the multichannel decomposition captures the global structure
    of the phase space,
    \item the VEGAS adaptation refines the sampling locally.
\end{itemize}

This mirrors the approach used in Munich, where both techniques are used
together to achieve stable and efficient integration of NLO contributions.

\subsection{Hadronic convolution}

The hadronic cross section is obtained as
\begin{align}
\sigma =
\sum_{a,b} \int dx_1 dx_2 \,
f_a(x_1,\mu_F) f_b(x_2,\mu_F)
\left[
\int d\Phi_{m+1} \, \mathcal{R}_{\text{finite}}
+
\int d\Phi_m \, \mathcal{V}_{\text{finite}}
\right].
\end{align}

\subsection{Numerical stability}

The combination of local subtraction, multichannel sampling,
and adaptive VEGAS integration ensures that:
\begin{itemize}
    \item the integrand remains finite across the entire phase space,
    \item singular regions are efficiently sampled,
    \item the Monte Carlo integration converges reliably.
\end{itemize}